\SourcePage{288}
\SGAsection{5}{与铅笔相伴的态射 $p:\widetilde X\to\mathbf P^1$ 的上同调研究}
\SGAsection{5.1}{限制同态 $H^*(\widetilde X)\to H^*(Y)$}
记 $Y$ 为 $p:\widetilde X\to\mathbf P^1$ 的一条\Term{smooth-fibre}{光滑纤维}，设 $Y=p^{-1}(t_0)=X\cdot H_{t_0}$。有如下\Term{commutative-diagram}{交换图}：
\[
\begin{tikzcd}[column sep=large,row sep=large]
\mathbf P^1 & \widetilde X\arrow[l,"p"']\arrow[d,"f"] & \widetilde\Delta\arrow[l,hook,"i"']\arrow[d,"g"]\arrow[r,"\sim"] & \Delta\times\mathbf P^1\\
\operatorname{Spec}(k)=\{t_0\}\arrow[u] & X & \Delta\arrow[l,hook,"j"]\arrow[ur,hook,"{\ell=\operatorname{id}\times\{t_0\}}"']\\
& Y\arrow[ul]\arrow[uu,hook,"k",bend left=20]\arrow[u,hook,"m"] & \Delta\arrow[l,hook,"h"']\arrow[u,equal]
\end{tikzcd}
\]
\SourcePage{289}
\begin{SGAresult}{命题}{5.1.1}
在 (3.1) 的假设下，设 $Y$ 是一条\Term{smooth-fibre}{光滑纤维}。
\textbf{(5.1.2)} \Term{restriction-map}{限制同态}
\[k^*:H^*(\widetilde X,\mathbf Z_\ell)\longrightarrow H^*(Y,\mathbf Z_\ell)\]
借助 (4.2.2) 可写为
\[H^*(X,\mathbf Z_\ell)\oplus H^{*-2}(\Delta,\mathbf Z_\ell(-1))\xrightarrow{m^*+h_*}H^*(Y,\mathbf Z_\ell).\]
\textbf{(5.1.3)} \Term{gysin-homomorphism}{Gysin 映射}
\[k_*:H^{*-2}(Y,\mathbf Z_\ell(-1))\longrightarrow H^*(\widetilde X,\mathbf Z_\ell)\]
可写为
\[H^{*-2}(Y,\mathbf Z_\ell(-1))\xrightarrow{m_*\oplus(-h^*)}H^*(X,\mathbf Z_\ell)\oplus H^{*-2}(\Delta,\mathbf Z_\ell(-1)).\]
\end{SGAresult}
\begin{proof}
$\widetilde X$ 的两个\Term{smooth-subscheme}{光滑子概形} $\widetilde\Delta$ 与 $Y$ 沿 $\Delta$ \Term{transverse-intersection}{横截相交}；换言之，下面的\Term{smooth-scheme}{光滑概形}方块是\Term{cartesian-square}{笛卡尔方块}：
\begin{equation}\tag{5.1.4}
\begin{tikzcd}
\Delta\arrow[r,hook,"\ell"]\arrow[d,hook,"h"']&\widetilde\Delta\arrow[d,hook,"i"]\\
Y\arrow[r,hook,"k"']&\widetilde X.
\end{tikzcd}
\end{equation}
因此有公式 [SGA 5，IV\SGAsourcegap{原 PDF 第289页引注“[SGA 5, IV”与右\Term{lie-bracket}{括号}之间留白，缺少具体定位。}]
\begin{equation}\tag{5.1.5}h_*\ell^*=k^*i_*,\qquad \ell_*h^*=i^*k_* .\end{equation}
\SourcePage{290}
复合
\[H^*(X,\mathbf Z_\ell)\oplus H^{*-2}(\Delta,\mathbf Z_\ell(-1))\xrightarrow{f^*+i_*g^*}H^*(\widetilde X,\mathbf Z_\ell)\xrightarrow{k^*}H^*(Y,\mathbf Z_\ell)\]
等于 $k^*f^*=m^*$ 与 $k^*i_*g^*=h_*\ell^*g^*$ 的和（由 (5.1.5)）；再由 $g\ell=\operatorname{id}_\Delta$，即得 (5.1.2)。复合
\[H^{*-2}(Y,\mathbf Z_\ell(-1))\xrightarrow{k_*}H^*(\widetilde X,\mathbf Z_\ell)\xrightarrow{f_*\oplus(-g_*i^*)}H^*(X,\mathbf Z_\ell)\oplus H^{*-2}(\Delta,\mathbf Z_\ell(-1))\]
的两个\Term{set-component}{分量}为 $f_*k_*=m_*$ 及 $-g_*i^*k_*=-g_*\ell_*h^*=-h^*$，从而得到 (5.1.3)。
\end{proof}
\begin{SGAresult}{推论}{5.1.6}
在 5.1.1 的假设下，当 $i\leq n-1$ 时，
\[m^*:H^i(X,\mathbf Z_\ell)\longrightarrow H^i(Y,\mathbf Z_\ell),\qquad k^*:H^i(\widetilde X,\mathbf Z_\ell)\longrightarrow H^i(Y,\mathbf Z_\ell)\]
的像相同。
\end{SGAresult}
\begin{proof}
由关于\Term{hyperplane-section}{超平面截面}的“\Term{weak-lefschetz-theorem}{弱 Lefschetz 定理}”（SGA 5 VII 7.1），当 $i\leq n-1$ 时，
\begin{equation}\tag{5.1.7}H^{i-2}(X,\mathbf Z_\ell(-1))\xrightarrow{j^*}H^{i-2}(\Delta,\mathbf Z_\ell(-1))\end{equation}
是\Term{isomorphism}{同构}。因而，当 $i\leq n-1$ 时，由 (5.1.2)，$k^*$ 的像等于 $m^*$ 的像与复合
\[H^{i-2}(X,\mathbf Z_\ell(-1))\xrightarrow{j^*}H^{i-2}(\Delta,\mathbf Z_\ell(-1))\xrightarrow{h_*}H^i(Y,\mathbf Z_\ell)\]
的像之和。
\SourcePage{291}
而由\Term{projection-formula}{投影公式}，
\[h_*j^*(x)=h_*\bigl(h^*(m^*(x))\bigr)=m^*(x)h_*(1)=\xi_Ym^*(x)=m^*(x\xi_X),\]
最后一步用了 $m^*(\xi_X)=\xi_Y$。所以 $\operatorname{Im}k^i\subset\operatorname{Im}m^i$，证毕。
\end{proof}
\SGAsection{5.2}{“Lefschetz 牛定理”（LV）}
\textbf{5.2.1。} 设 $S$ 是\Term{algebraically-closed-field}{代数闭域} $k$ 上\Term{dimension}{维数}为 $n$ 的\Term{projective-scheme}{射影概形}，并赋予\Term{projective-embedding}{射影嵌入} $S\hookrightarrow\mathbf P$。记
\[u\in H^2(S,\mathbf Z_\ell(1))\]
为“\Term{hyperplane-section}{超平面截面}的\Term{cohomology-class}{上同调类}”（也就是说，$u$ 是 $\mathscr O_{\mathbf P}(1)$ 的\Term{first-chern-class}{第一 Chern 类}沿\Term{embedding}{嵌入} $S\hookrightarrow\mathbf P$ 的\Term{inverse-image}{逆像}；后者属于 $H^2(\mathbf P,\mathbf Z_\ell(1))$）。记
\[L_S:H^*(S,\mathbf Q_\ell(j))\longrightarrow H^{*+2}(S,\mathbf Q_\ell(j+1))\]
为“乘以 $u$”这一运算。
\textbf{5.2.2。} 如果对每个（或某个）整数 $j$ 及每个整数 $i>0$，映射
\[L_S^i:H^{n-i}(S,\mathbf Q_\ell(j))\longrightarrow H^{n+i}(S,\mathbf Q_\ell(i+j))\]
都是\Term{isomorphism}{同构}，就说 $S$（连同 $u$）满足 (LV)，即“\Term{lefschetz-cow-theorem}{Lefschetz 牛定理}”（théorème de Lefschetz vache）。
\SourcePage{292}
\textbf{5.2.2.1。} 猜想当 $S$ \Term{smooth}{光滑}时，此条件成立。当 $k$ 的\Term{characteristic}{特征}为零时，这是成立的 [6]；由“\Term{specialization-theorem}{专化定理}”SGA 4 XVI 2.2 可\Term{pushout}{推出}，如果 $S$ 连同其\Term{projective-embedding}{射影嵌入}可以\Term{lifting}{提升}到\Term{characteristic-zero}{特征零}，结论同样成立。因此，当 $S$ 是\Term{smooth-complete-intersection}{光滑完全交}时，此条件成立。还知道，当 $S$ \Term{smooth}{光滑}且\Term{dimension}{维数} $\leq2$ 时 [2]，以及当 $S$ 是\Term{abelian-variety}{阿贝尔簇}时，结论成立（后一情形既可“直接”证明，见 LIEBERMAN [4] 及 [2]；也可由 MUMFORD 的\Term{theorem}{定理} [5] \Term{pushout}{推出}，该\Term{theorem}{定理}说\Term{polarized-abelian-variety}{极化阿贝尔簇}可连同其\Term{polarization}{极化}一起\Term{lifting}{提升}到\Term{characteristic-zero}{特征零}）。回顾一下，对\Term{projective-variety}{射影簇} $S$ 及 $i\leq n$，\Term{cohomology}{上同调}的“\Term{primitive-part}{本原部分}”定义为
\begin{equation}\tag{5.2.2.2}
\operatorname{Prim}^i(S,\mathbf Q_\ell(j))=\operatorname{Ker}\bigl(L^{n-i+1}:H^i(S,\mathbf Q_\ell(j))\to H^{2n-i+2}(S,\mathbf Q_\ell(j+2n-2i+2))\bigr).
\end{equation}
如果 $S$ 满足 (LV)，则当 $i\leq n$ 时有\Term{decomposition}{分解}
\begin{equation}\tag{5.2.2.3}
H^i(S,\mathbf Q_\ell(j))\xleftarrow[\sim]{1+L}\operatorname{Prim}^i(S,\mathbf Q_\ell(j))\oplus H^{i-2}(S,\mathbf Q_\ell(j-1)).
\end{equation}
\begin{SGAresult}{命题}{5.2.3}
设 $X$ \Term{projective}{射影}、\Term{smooth}{光滑}且满足 (LV)。以 $\tau:Y\hookrightarrow X$ 表示一个\Term{smooth-hyperplane-section}{光滑超平面截面}的\Term{inclusion}{包含映射}。则
\begin{equation}\tag{5.2.3.1}\operatorname{Prim}^n(X,\mathbf Q_\ell(j))=\operatorname{Ker}\bigl(\tau^*:H^n(X,\mathbf Q_\ell(j))\to H^n(Y,\mathbf Q_\ell(j))\bigr).\end{equation}
\end{SGAresult}
\begin{proof}
由 $m:Y\to X$ 的\Term{projection-formula}{投影公式}，有\Term{commutative-diagram}{交换图}
\SourcePage{293}
\begin{equation}\tag{5.2.3.2}
\begin{tikzcd}
H^n(X,\mathbf Q_\ell(j))\arrow[rr,"L"]\arrow[dr,"\tau^*"']&&H^{n+2}(X,\mathbf Q_\ell(j+1))\\
&H^n(Y,\mathbf Q_\ell(j))\arrow[ur,"\tau_*"']&
\end{tikzcd}
\end{equation}
由\Term{duality}{对偶性}及\Term{weak-lefschetz-theorem}{弱 Lefschetz 定理}，
\[\tau_*:H^n(Y,\mathbf Q_\ell(j))\longrightarrow H^{n+2}(X,\mathbf Q_\ell(j+1))\]
是一个\Term{isomorphism}{同构}的\Term{transpose}{转置}，因此是\Term{isomorphism}{同构}，结论随即得到。
\end{proof}
\begin{SGAresult}{\Term{lemma}{引理}}{5.2.4}
设 $X$ \Term{projective}{射影}且\Term{smooth}{光滑}，$\tau:Y\hookrightarrow X$ 是一个\Term{smooth-hyperplane-section}{光滑超平面截面}的\Term{inclusion}{包含映射}。则 $X$ 满足 (LV) 当且仅当下列两个条件成立。
\textbf{5.2.4.1。} $Y$ 满足 (LV)。
\textbf{5.2.4.2。} $H^{n-1}(Y,\mathbf Q_\ell)$ 上的\Term{cup-product}{杯积}\Term{restriction}{限制}到\Term{vector-subspace}{子空间} $\operatorname{Im}H^{n-1}(X,\mathbf Q_\ell)$ 后，是\Term{nondegenerate-form}{非退化形式}。
\end{SGAresult}
\begin{proof}
有
\begin{equation}\tag{5.2.4.3}\left\{\begin{aligned}L_X&=\tau_*\tau^*,\\L_Y&=\tau^*\tau_*\qquad\text{（由 4.2.11）},\end{aligned}\right.\end{equation}
\SourcePage{294}
因而
\begin{equation}\tag{5.2.4.4}L_X^{i+1}=\tau_*L_Y^i\tau^*.\end{equation}
在\Term{commutative-diagram}{交换图}
\begin{equation}\tag{5.2.4.5}
\begin{tikzcd}[column sep=large]
H^{n-i-1}(X,\mathbf Q_\ell(j))\arrow[r,"L_X^{i+1}"]\arrow[d,"\tau^*"']&H^{n+i+1}(X,\mathbf Q_\ell(j+i+1))\\
H^{n-1-i}(Y,\mathbf Q_\ell(j))\arrow[r,"L_Y^i"']&H^{n-1+i}(Y,\mathbf Q_\ell(j+i))\arrow[u,"\tau_*"']
\end{tikzcd}
\end{equation}
中，由“\Term{weak-lefschetz-theorem}{弱 Lefschetz 定理}”，$i\geq1$ 时 $\tau^*$ 是\Term{isomorphism}{同构}；由\Term{duality}{对偶性}，$\tau_*$ 也是\Term{isomorphism}{同构}。所以 $X$ 满足 (LV) 当且仅当 $Y$ 满足 (LV)，且映射
\begin{equation}\tag{5.2.4.6}L_X:H^{n-1}(X,\mathbf Q_\ell(j))\longrightarrow H^{n+1}(X,\mathbf Q_\ell(j+1))\end{equation}
是\Term{isomorphism}{同构}。要完成证明，只需说明 5.2.4.2 与 5.2.4.6 \Term{equivalence}{等价}。为此，设 $a,b\in H^{n-1}(X,\mathbf Q_\ell(j))$。由 $X$ 上的\Term{poincare-duality}{Poincaré 对偶性}，(5.2.4.6) \Term{equivalence}{等价}于：
\textbf{(5.2.4.7)} \Term{bilinear-form}{双线性形式}
\[H^{n-1}(X,\mathbf Q_\ell(j))\times H^{n-1}(X,\mathbf Q_\ell(j))\longrightarrow H^{2n}(X,\mathbf Q_\ell(2j+2)),\qquad(a,b)\longmapsto a\cdot L_X(b)\]
是\Term{nondegenerate-form}{非退化的}。
\SourcePage{295}
但由 (5.2.5) 及\Term{projection-formula}{投影公式}，有
\[a\cdot L_X(b)=a\cdot\tau_*(\tau^*(b))=\tau_*\bigl(\tau^*(a)\cdot\tau^*(b)\bigr),\]
因此 (5.2.4.7) 与 (5.2.4.2) \Term{equivalence}{等价}。
\end{proof}
\begin{SGAresult}{推论}{5.2.5}
设 $X$ \Term{projective}{射影}、\Term{smooth}{光滑}且\Term{dimension}{维数}为 $n$，$\tau:Y\hookrightarrow X$ 是一个\Term{smooth-hyperplane-section}{光滑超平面截面}的\Term{inclusion}{包含映射}。假设 $X$ 满足条件 (LV)。则当 $i\geq n$ 时，映射
\[\tau^*:H^i(X,\mathbf Q_\ell(j))\longrightarrow H^i(Y,\mathbf Q_\ell(j))\]
是\Term{surjectivity}{满射}。
\end{SGAresult}
\begin{proof}
由 (5.2.4.3)，对 $i\geq0$，有 $L_Y^{i+1}=\tau^*L_X^i\tau_*$，从而有\Term{commutative-diagram}{交换图}
\[
\begin{tikzcd}[column sep=large]
H^{n-i}(X,\mathbf Q_\ell(j+1))\arrow[r,"L_X^i"]&H^{n-i}(X,\mathbf Q_\ell(j+i+1))\arrow[d,"\tau^*"]\\
H^{n-i-2}(Y,\mathbf Q_\ell(j))\arrow[u,"\tau_*"]\arrow[r,"L_Y^{i+1}"',"\sim"]&H^{n+i}(Y,\mathbf Q_\ell(j+i+1)),
\end{tikzcd}
\]
这就给出所需结论，证毕。
\end{proof}
\SourcePage{296}
\begin{SGAresult}{\Term{lemma}{引理}}{5.2.6}
设 $f:F\to B$ 是\Term{proper-morphism}{固有}且\Term{smooth-morphism}{光滑的态射}，$B$ \Term{connected}{连通}。记 $\mathscr O_F(1)$ 为 $F$ 上的一个\Term{relatively-ample}{相对丰沛模}，并给每条\Term{geometric-fibre}{几何纤维} $F_{\bar b}$ 赋予由 $\mathscr O_F(1)$ 诱导的\Term{ample}{丰沛模}。如果其中一条\Term{geometric-fibre}{几何纤维}满足条件 (LV)，则其余各条也都满足。
\end{SGAresult}
\begin{proof}
记 $\xi_F\in H^2(F,\mathbf Z_\ell(1))$ 为 $\mathscr O_F(1)$ 的\Term{first-chern-class}{第一 Chern 类}。此类及其幂 $\xi_F^i\in H^{2i}(F,\mathbf Z_\ell(i))$ 通过\Term{cup-product}{杯积}诱导映射
\begin{equation}\tag{5.2.6.1}\xi_F^i:R^{n-i}f_*\mathbf Q_\ell(j)\longrightarrow R^{n+i}f_*\mathbf Q_\ell(j+i),\end{equation}
其中 $n$ 表示\Term{fibre}{纤维}的\Term{relative-dimension}{相对维数}。由\Term{proper-base-change}{固有态射的基变换定理} [SGA 4 XII 5.2]，所有\Term{geometric-fibre}{几何纤维}满足条件 (LV)，当且仅当 (5.2.6.1) 中的全部映射都是\Term{isomorphism}{同构}。而由\Term{specialization-theorem}{专化定理} [SGA 4 XVI 2.2]，\Term{sheaf}{层} $R^if_*\mathbf Q_\ell(j)$ 是\Term{twisted-constant}{扭常值的}，因此 (5.2.6.1) 中的映射是\Term{isomorphism}{同构}，当且仅当它们在给定的一个\Term{geometric-point}{几何点}上是\Term{isomorphism}{同构}，证毕。
\end{proof}
\begin{SGAresult}{推论}{5.2.7}
设 $X$ \Term{projective}{射影}且\Term{smooth}{光滑}。若一个\Term{smooth-hyperplane-section}{光滑超平面截面}满足 (LV)，则每个\Term{smooth-hyperplane-section}{光滑超平面截面}都满足 (LV)。
\end{SGAresult}
\begin{proof}
事实上，$X$ 的\Term{smooth-hyperplane-section}{光滑超平面截面}的\Term{universal-family}{泛族}满足 5.2.5 的假设。
\end{proof}
\SourcePage{297}
\SGAsection{5.3}{条件 (A)}
\textbf{5.3.1。} 设
\[p:\widetilde X\longrightarrow\mathbf P^1\]
是与一个\Term{pencil}{铅笔}相伴的“\Term{fibration}{纤维化}”；其\Term{pencil-axis}{轴}与 $X$ \Term{transverse-intersection}{横截相交}，并假设其\Term{generic-fibre}{一般纤维}是\Term{smooth}{光滑的}。
记 $U$ 为 $\mathbf P^1$ 的一个\Term{nonempty}{非空}\Term{open-subset}{开集}，使得 $p$ 在其上\Term{smooth-morphism}{光滑}，并记\Term{inclusion}{包含映射}为
\begin{equation}\tag{5.3.1.1}\nu:U\hookrightarrow\mathbf P^1.\end{equation}
由于\Term{sheaf}{层} $R^ip_*\mathbf Q_\ell(j)$ 在 $U$ 上\Term{twisted-constant}{扭常值}，可见\Term{sheaf}{层}
\[\nu_*\nu^*R^ip_*\mathbf Q_\ell(j)\]
本质上与 $U$ 的选择无关，且\Term{canonical}{典范}“\Term{adjunction}{伴随}”\Term{homomorphism}{同态}
\begin{equation}\tag{5.3.2}R^ip_*\mathbf Q_\ell(j)\longrightarrow\nu_*\nu^*R^ip_*\mathbf Q_\ell(j)\end{equation}
也与 $U$ 的选择无关。
\textbf{5.3.3。} 注意，对任意\Term{hyperplane-pencil}{超平面截面铅笔}，\Term{weak-lefschetz-theorem}{弱 Lefschetz 定理}说明，当 $i\leq n-2$ 时，\Term{immersion}{浸入}
$\widetilde X\xhookrightarrow{f\times p}X\times\mathbf P^1$ 给出 $\mathbf P^1$ 上的\Term{sheaf}{层}\Term{isomorphism}{同构}
\[R^ip_*\mathbf Q_\ell(j)\xleftarrow{\sim}H^i(X,\mathbf Q_\ell(j))_{\mathbf P^1},\qquad i\leq n-2;\]
\SourcePage{298}
也就是说，$i\leq n-2$ 时 $R^ip_*\mathbf Q_\ell(j)$ 是\Term{constant-sheaf}{常值层}；尤其 (5.3.2) 中的映射在 $i\leq n-2$ 时是\Term{isomorphism}{同构}。
\textbf{5.3.4。} 由\Term{vanishing-cycles}{消失圈}的\Term{local-scheme}{局部}理论，特别是\Term{picard-lefschetz-formula}{Picard--Lefschetz 公式}（第 XV 讲 3.4），知道对一个\Term{lefschetz-pencil}{Lefschetz 铅笔}，\Term{homomorphism}{同态} 5.3.2 在 $i\neq n$ 时是\Term{isomorphism}{同构}，在 $i=n$ 时是\Term{surjectivity}{满射}。
\textbf{5.3.5。} 如果 (5.3.2) 中的所有\Term{homomorphism}{同态}都是\Term{isomorphism}{同构}（\Term{equivalence}{等价}地，对所有 $i$ 及 $j=0$ 都是\Term{isomorphism}{同构}），就说 $p$ 或产生它的\Term{pencil}{铅笔}满足条件 (A)。下文（6.3、6.4）将会看到，这个条件 (A) 颇有\Term{universal}{普遍}成立的趋势。
\begin{SGAresult}{\Term{lemma}{引理}}{5.4}
如果 $X$ 的某个\Term{smooth-hyperplane-section}{光滑超平面截面}（由 5.2.7，也就是每个）满足 (LV)（例如当 $X$ 满足 (LV) 时，见 (5.2.4)），且 $p$ 满足假设 (A)，则每条\Term{fibre}{纤维}（包括每条\Term{singular-fibre}{奇异纤维}）都满足 (LV)，而且 $i\neq n-1$ 时，\Term{sheaf}{层} $R^ip_*\mathbf Q_\ell(j)$ 是\Term{constant-sheaf}{常值的}。
\end{SGAresult}
\begin{proof}
以
\[L:R^ip_*\mathbf Q_\ell(j)\longrightarrow R^{i+2}p_*\mathbf Q_\ell(j+1)\]
表示乘以 $f^*c_1(\mathscr O_X(1))\in H^2(\widetilde X,\mathbf Z_\ell(1))$ 的\Term{morphism}{态射}；逐\Term{fibre}{纤维}来看，这就是 (5.2.1) 所考虑的运算。由假设，在 $U$ 上有\Term{isomorphism}{同构}
\SourcePage{299}
\begin{equation}\tag{5.4.1}\nu^*(L^i):\nu^*R^{n-1-i}p_*\mathbf Q_\ell(j)\xrightarrow{\sim}\nu^*R^{n-1+i}p_*\mathbf Q_\ell(j+i),\end{equation}
对 $i>0$ 及所有 $j$ 都成立。应用\Term{functor}{函子} $\nu_*$，得到\Term{isomorphism}{同构}
\begin{equation}\tag{5.4.2}(\nu_*\nu^*)L^i:\nu_*\nu^*R^{n-1-i}p_*\mathbf Q_\ell(j)\xrightarrow{\sim}\nu_*\nu^*R^{n-1+i}p_*\mathbf Q_\ell(j+i);\end{equation}
由 (A)，它们可与\Term{morphism}{态射}
\begin{equation}\tag{5.4.3}L^i:R^{n-1-i}p_*\mathbf Q_\ell(j)\xrightarrow{\sim}R^{n-1+i}p_*\mathbf Q_\ell(j+i)\end{equation}
等同。再由 5.3.3 即得结论。
\end{proof}
\textbf{5.4.4。} 记 $Y$ 为 $p:\widetilde X\to\mathbf P^1$ 的一条\Term{smooth-fibre}{光滑纤维}。由“\Term{weak-lefschetz-theorem}{弱 Lefschetz 定理}”，\Term{restriction-map}{限制映射}
\[H^{n-1}(X,\mathbf Q_\ell(j))\xhookrightarrow{m^*}H^{n-1}(Y,\mathbf Q_\ell(j))\]
是\Term{injectivity}{单射}。已经看到（5.2.4），如果 $X$ 满足 (LV)，则 $H^{n-1}(Y,\mathbf Q_\ell(j))$ 上的\Term{cup-product}{杯积}\Term{restriction}{限制}到 $m^*H^{n-1}(X,\mathbf Q_\ell(j))$ 后仍\Term{nondegenerate-form}{非退化}。因此 $H^{n-1}(X,\mathbf Q_\ell(j))$ 在 $H^{n-1}(Y,\mathbf Q_\ell(j))$ 中的\Term{orthogonal-complement}{正交补}是它的一个\Term{complementary-subspace}{补空间}，依照经典术语，称为 $Y$ 的\Term{vanishing-cohomology}{消失上同调空间}。在 $\nu_*R^{n-1}p_*\mathbf Q_\ell(j)$ 中定义相应的\Term{subsheaf}{子层}（这里 $\nu:U\to\mathbf P^1$ 是 $p$ 的一个\Term{nonempty}{非空}\Term{smooth-locus}{光滑开集}），即 $\mathbf P^1$ 上的“\Term{vanishing-cohomology-sheaf}{消失上同调层}” $E^{n-1}p_*\mathbf Q_\ell(j)$，如下：
\SourcePage{300}
\begin{equation}\tag{5.4.4.1}E^{n-1}p_*\mathbf Q_\ell(j)=\nu_*\left\{\begin{gathered}\nu^*R^{n-1}p_*\mathbf Q_\ell(j)\text{ 中}\text{常值子层 }H^{n-1}(X,\mathbf Q_\ell(j))_U\text{ 的正交补}\end{gathered}\right\}.\end{equation}
显然，此\Term{sheaf}{层}与 $U$ 的选择无关。
\begin{SGAresult}{\Term{lemma}{引理}}{5.5}
如果 $X$ 满足 (LV)，且 (5.3.2) 中的映射在 $i=n-1$ 时是\Term{isomorphism}{同构}，例如 $p$ 来自一个\Term{lefschetz-pencil}{Lefschetz 铅笔}时（见 (5.3.4)），则有\Term{direct-sum-decomposition}{直和分解}
\begin{equation}\tag{5.5.1}R^{n-1}p_*\mathbf Q_\ell(j)\simeq H^{n-1}(X,\mathbf Q_\ell(j))_{\mathbf P^1}\oplus E^{n-1}p_*\mathbf Q_\ell(j).\end{equation}
\end{SGAresult}
\begin{proof}
由 (LV)，在 $U$ 上有这样的\Term{decomposition}{分解}：
\begin{equation}\tag{5.5.2}\nu^*R^{n-1}p_*\mathbf Q_\ell(j)\simeq H^{n-1}(X,\mathbf Q_\ell(j))_U\oplus\nu^*E^{n-1}p_*\mathbf Q_\ell(j);\end{equation}
对它应用 $\nu_*$ 即可。
\end{proof}
\begin{SGAresult}{\Term{theorem}{定理}}{5.6}
设 $X$ 是满足 (LV) (5.2.2) 的\Term{smooth}{光滑}\Term{irreducible}{不可约}\Term{projective-variety}{射影簇}。设 $\{H_\tau\}$ 是一个\Term{hyperplane-pencil}{超平面铅笔}，其\Term{generic-member}{一般成员}及其\Term{pencil-axis}{轴}分别沿 $X_\eta$ 及 $\Delta$ 与 $X$ \Term{transverse-intersection}{横截相交}。假设\Term{projection}{投影} $p:\widetilde X\to\mathbf P^1$ 满足条件 (5.3.2)。则：
\textbf{(5.6.1)} \Term{leray-spectral-sequence}{Leray 谱序列}
\[E_2^{p,q}=H^p(\mathbf P^1,R^qp_*\mathbf Q_\ell)\Longrightarrow H^{p+q}(\widetilde X,\mathbf Q_\ell)\]
\Term{degenerate-spectral-sequence}{退化}。（注意：如果不假设 $X$ 满足条件 (LV)，而仅假设某个\Term{smooth-hyperplane-section}{光滑超平面截面} $Y$ 满足，结论仍然成立。）
\SourcePage{301}
\textbf{5.6.2。} 以 $\eta$ 表示 $\mathbf P^1$ 的\Term{generic-point}{泛点}，$\bar\eta$ 表示\Term{geometric-generic-point}{几何泛点}，$X_{\bar\eta}$ 表示 $p$ 的\Term{geometric-generic-fibre}{几何一般纤维}，$\pi$ 表示 $k(\bar\eta)/k(\eta)$ 的\Term{galois-group}{Galois 群}。当 $q\neq n-1$ 时，有\Term{canonical}{典范}\Term{isomorphism}{同构}：
\begin{align}E_2^{0,q}\;(=H^0(\mathbf P^1,R^qp_*\mathbf Q_\ell))&\simeq H^q(X_{\bar\eta},\mathbf Q_\ell),\tag{5.6.2.1}\\
E_2^{2,q}\;(=H^2(\mathbf P^1,R^qp_*\mathbf Q_\ell))&\simeq H^q(X_{\bar\eta},\mathbf Q_\ell(-1)),\tag{5.6.2.2}\\
E_2^{p,q}\;(=H^p(\mathbf P^1,R^qp_*\mathbf Q_\ell))&=0,\quad p\neq0,2.\tag{5.6.2.3}\end{align}
当 $q=n-1$ 时，有
\begin{align}
E_2^{0,n-1}\;(=H^0(\mathbf P^1,R^{n-1}p_*\mathbf Q_\ell))&\simeq H^{n-1}(X_{\bar\eta},\mathbf Q_\ell)^\pi\simeq H^{n-1}(X,\mathbf Q_\ell),\tag{5.6.2.4}\\
E_2^{2,n-1}\;(=H^2(\mathbf P^1,R^{n-1}p_*\mathbf Q_\ell))&\simeq H^{n-1}(X_{\bar\eta},\mathbf Q_\ell(-1))^\pi\simeq H^{n-1}(X,\mathbf Q_\ell(-1)).\tag{5.6.2.5}
\end{align}
（尚未列出的 $E_2^{1,n-1}$ 将在下文 5.7 中求出。）
\end{SGAresult}
\begin{proof}
由 (5.4.3)，可以应用\Term{theorem}{定理} 1.1，得到 (5.6.1)。至于 (5.6.2) 的前三个公式，它们来自 $q\neq n-1$ 时 $R^qp_*\mathbf Q_\ell$ 为\Term{constant-sheaf}{常值层}这一事实（5.4）；对 (5.6.2.3)，还须使用 $\mathbf P^1$ \Term{simply-connected}{单连通}这一点。(5.6.2.4) 的第一个等式来自
\SourcePage{302}
$R^{n-1}p_*\mathbf Q_\ell\xrightarrow{\sim}\nu_*\nu^*R^{n-1}p_*\mathbf Q_\ell$，因而
$H^0(\mathbf P^1,R^{n-1}p_*\mathbf Q_\ell)\simeq H^0(U,\nu^*R^{n-1}p_*\mathbf Q_\ell)$；又因为 $\nu^*R^{n-1}p_*\mathbf Q_\ell$ 是 $U$ 上的“\Term{twisted-constant}{扭常值}”\Term{sheaf}{层}，所以\Term{equivalence}{等价}于 $\pi$-模 $H^{n-1}(X_{\bar\eta},\mathbf Q_\ell)$。(5.6.2.4) 的第二个等式来自\Term{degenerate-spectral-sequence}{退化}，并用到 5.1.6。
再看 (5.6.2.5)。记 $F^*$ 为\Term{spectral-sequence}{谱序列}在 $H^*(\widetilde X,\mathbf Q_\ell)$ 上定义的\Term{filtration}{滤过}。由\Term{degenerate-spectral-sequence}{退化}，
$F^2H^{n+1}(\widetilde X,\mathbf Q_\ell)=E_2^{2,n-1}$，而 $\operatorname{gr}^1H^{n+1}(\widetilde X,\mathbf Q_\ell)=E_2^{1,n}=0$（5.6.2.3），故有（见 5.1.0）
\begin{equation}\tag{5.6.3}\begin{aligned}
E_2^{2,n-1}&=F^1H^{n+1}(\widetilde X,\mathbf Q_\ell)\\
&=\operatorname{Ker}\bigl(H^{n+1}(\widetilde X,\mathbf Q_\ell)\to H^0(\mathbf P^1,R^{n+1}p_*\mathbf Q_\ell)\bigr)\\
&\overset{(A)}=\operatorname{Ker}\bigl(k^*:H^{n+1}(\widetilde X,\mathbf Q_\ell)\to H^{n+1}(X_{\bar\eta},\mathbf Q_\ell)\bigr).
\end{aligned}\end{equation}
由 (5.1.2)，$k^*$ 可改写为
\begin{equation}\tag{5.6.4}H^{n+1}(X,\mathbf Q_\ell)\oplus H^{n-1}(\Delta,\mathbf Q_\ell(-1))\xrightarrow{m^*+h_*}H^{n+1}(X_{\bar\eta},\mathbf Q_\ell),\end{equation}
再由 (5.2.2.3)，又可改写为
\begin{equation}\tag{5.6.5}\begin{gathered}
\operatorname{Prim}^{n-1}(X,\mathbf Q_\ell(-1))\oplus H^{n-3}(X,\mathbf Q_\ell(-2))\oplus H^{n-1}(\Delta,\mathbf Q_\ell(-1))\\
\xrightarrow{m^*L+m^*L^2+h_*}H^{n+1}(X_{\bar\eta},\mathbf Q_\ell).
\end{gathered}\end{equation}
而映射
\[H^{n-3}(X,\mathbf Q_\ell(-2))\xrightarrow{m^*L^2}H^{n+1}(X_{\bar\eta},\mathbf Q_\ell)\]
由 (LV) 是\Term{isomorphism}{同构}，因为 $m^*L_X^2=L_{X_{\bar\eta}}^2\cdot m^*$（见 (5.2.4) 及 (5.2.5)）。因而 (5.6.8) 的核经 $\operatorname{pr}_1\oplus\operatorname{pr}_3$ 可等同于
\SourcePage{303}
$\operatorname{Prim}^{n-1}(X,\mathbf Q_\ell(-1))\oplus H^{n-1}(\Delta,\mathbf Q_\ell(-1))$。由 (LV)，
\[H^{n-1}(\Delta,\mathbf Q_\ell(-1))\xleftarrow[\sim]{L}H^{n-3}(\Delta,\mathbf Q_\ell(-2))\xleftarrow{\sim}H^{n-3}(X,\mathbf Q_\ell(-2)),\]
最后一个\Term{isomorphism}{同构}来自\Term{weak-lefschetz-theorem}{弱 Lefschetz 定理}。因此由 (5.2.2.3)，有
\begin{equation}\tag{5.6.6}E_2^{2,n-1}\simeq H^{n-1}(X,\mathbf Q_\ell(-1)),\end{equation}
而 (5.6.2.5) 的另一个等式可由对 (5.6.2.4) 作“\Term{tate-twist}{扭转}”得到。
\end{proof}
\textbf{注记 5.6.7。} DELIGNE 向我解释，可以在明显更弱的假设下证明 (5.6.1)：
\begin{SGAresult}{\Term{theorem}{定理}}{5.6.8}
设 $X$ 在\Term{algebraically-closed-field}{代数闭域} $k$ 上\Term{smooth}{光滑}，\Term{dimension}{维数}为 $n$，且
\[p:X\longrightarrow\mathbf P^1_k\]
是一个\Term{projective-morphism}{射影态射}，其\Term{generic-fibre}{一般纤维}是\Term{smooth}{光滑的}。假设：
\textbf{(5.6.8.1)} \Term{geometric-generic-fibre}{几何一般纤维}，赋予\Term{relatively-ample}{相对丰沛模} $\mathscr O_X(1)$ 的\Term{restriction}{限制}，满足条件 (LV)；
\textbf{(5.6.8.2)} 映射 (5.3.2) 在 $i\neq n$ 时是\Term{isomorphism}{同构}，在 $i=n$ 时是\Term{surjectivity}{满射}（特别地，当 $p$ 来自一个\Term{lefschetz-pencil}{Lefschetz 铅笔}时，此条件成立，见 5.3.4）。则\Term{leray-spectral-sequence}{Leray 谱序列}
\[E_2^{p,q}:H^p(\mathbf P^1,R^qp_*\mathbf Q_\ell)\Longrightarrow H^{p+q}(X,\mathbf Q_\ell)\]
\Term{degenerate-spectral-sequence}{退化}。
\end{SGAresult}
\textbf{注记。} 把 $\mathbf P^1_k$ 换成任意\Term{smooth}{光滑}\Term{irreducible}{不可约}\Term{projective-curve}{射影曲线}，同一\Term{theorem}{定理}仍成立，证明也相同。
\SourcePage{304}
\begin{proof}
在证明中，将\Term{sheaf}{层} $\mathbf Q_\ell(1)$ 与 $\mathbf Q_\ell$ 等同（由于已假设 $k$ \Term{algebraically-closed}{代数闭}，这是可以的）。显然可假设 $X$ \Term{connected}{连通}且 $p$ \Term{surjective}{满射}，因此 $p$ \Term{flat}{平坦}，其\Term{fibre}{纤维}\Term{dimension}{维数}为 $n-1$。记
\[\xi_X\in H^2(X,\mathbf Q_\ell)\]
为 $\mathscr O_X(1)$ 的\Term{first-chern-class}{第一 Chern 类}。设 $\nu:U\hookrightarrow\mathbf P^1$ 是 $p$ 的一个\Term{smooth-locus}{光滑开集}的\Term{inclusion}{包含映射}。由 5.6.8.1 及 5.2.6，\Term{cup-product}{杯积}映射
\[\xi_X^i:\nu^*R^{n-1-i}p_*\mathbf Q_\ell\longrightarrow\nu^*R^{n-1+i}p_*\mathbf Q_\ell\]
在 $i\geq1$ 时是\Term{isomorphism}{同构}；因此应用\Term{functor}{函子} $\nu_*$，并用 5.6.8.2，当 $i\geq2$ 时得到\Term{isomorphism}{同构}
\begin{equation}\tag{5.6.8.3}\xi_X^i:R^{n-1-i}p_*\mathbf Q_\ell\xrightarrow{\sim}R^{n-1+i}p_*\mathbf Q_\ell,\end{equation}
当 $i=1$ 时得到\Term{isomorphism}{同构}
\begin{equation}\tag{5.6.8.4}\xi_X:R^{n-2}p_*\mathbf Q_\ell\xrightarrow{\sim}\nu_*\nu^*R^np_*\mathbf Q_\ell.\end{equation}
后一个映射\Term{decomposition}{分解}为
\begin{equation}\tag{5.6.8.5}R^{n-2}p_*\mathbf Q_\ell\xrightarrow{\xi_X}R^np_*\mathbf Q_\ell\xrightarrow{\text{伴随映射}}\nu_*\nu^*R^np_*\mathbf Q_\ell,\end{equation}
\SourcePage{305}
因而有\Term{canonical-decomposition}{典范分解}
\begin{equation}\tag{5.6.8.6}R^np_*\mathbf Q_\ell\xleftarrow[\sim]{\xi_X+\operatorname{ind}}R^{n-2}p_*\mathbf Q_\ell\oplus\operatorname{Local}^n\mathbf Q_\ell,\end{equation}
这里置
\begin{equation}\tag{5.6.8.7}\operatorname{Local}^n\mathbf Q_\ell=\operatorname{Ker}\bigl(R^np_*\mathbf Q_\ell\to\nu_*\nu^*R^np_*\mathbf Q_\ell\bigr).\end{equation}
注意，\Term{sheaf}{层} $\operatorname{Local}^n\mathbf Q_\ell$ 集中在 $\mathbf P^1$ 上 $p$ 的\Term{fibre}{纤维}不\Term{smooth}{光滑}的那些点；因而再用 (5.6.8.3)，得到
\begin{equation}\tag{5.6.8.8}H^j(\mathbf P^1,R^{n-2-i}p_*\mathbf Q_\ell)\xrightarrow[\sim]{\xi_X}H^j(\mathbf P^1,R^{n-1+i}p_*\mathbf Q_\ell),\qquad i,j\geq1.\end{equation}
现在证明\Term{degenerate-spectral-sequence}{退化性}。因为除 $p=0,1,2$ 外 $E_2^{p,q}=0$，只需证明运算
\[d_2:H^0(\mathbf P^1,R^ip_*\mathbf Q_\ell)\longrightarrow H^2(\mathbf P^1,R^{i-1}p_*\mathbf Q_\ell)\]
为零。对 $q\leq n-1$，置
\begin{equation}\tag{5.6.8.9}\operatorname{Prim}(E_2^{p,q})=\operatorname{Ker}\bigl(\xi_X^{n-q}:E_2^{p,q}\to E_2^{p,2n-q}\bigr).\end{equation}
由 (5.6.8.3) 及 (5.6.8.6)，当 $q\leq n-1$ 时，
\[\xi_X:H^p(\mathbf P^1,R^{n-2}p_*\mathbf Q_\ell)\longrightarrow H^p(\mathbf P^1,R^1p_*\mathbf Q_\ell)\]
为\Term{injectivity}{单射}，因此有
\begin{equation}\tag{5.6.8.10}E_2^{p,q}\xleftarrow[\sim]{1+\xi_X+\xi_X^2+\cdots}\operatorname{Prim}(E_2^{p,q})\oplus\operatorname{Prim}(E_2^{p,q-2})\oplus\operatorname{Prim}(E_2^{p,q-4}).\end{equation}
\SourcePage{306}
要说明 $q\leq n-1$ 时 $d_2^{p,q}=0$，只需注意\Term{commutative-diagram}{交换图}
\[
\begin{tikzcd}[column sep=large]
\operatorname{Prim}(E_2^{p,q})\arrow[r,"{\xi_X^{n-q}=0}"]\arrow[d,"{d_2^{p,q}}"']&E_2^{p,2n-q}\arrow[d,"{d_2^{p,2n-q}}"]\\
E_2^{p+2,q-1}\arrow[r,"\xi_X^{n-1-(q-1)}"',"\sim"]&E_2^{p+2,2n-q-1}
\end{tikzcd}
\]
（下面的映射由 (5.6.8.8) 是\Term{isomorphism}{同构}）表明：当 $q\leq n-1$ 时，$d_2^{p,q}$ 消去 $\operatorname{Prim}(E_2^{p,q})$，故由 (5.6.8.10)，$d_2$ 消去 $E_2^{p,q}$。由 (5.6.8.3)，这还表明 $q\geq n+1$ 时 $d_2$ 消去 $E_2^{p,q}$。因此只剩证明\Term{homomorphism}{同态}
\[d_2^{0,n}:E_2^{0,n}\longrightarrow E_2^{2,n-1}\]
为零；\Term{equivalence}{等价}地，证明
\begin{equation}\tag{5.6.8.12}\dim E_3^{2,n-1}=\dim E_2^{2,n-1}.\end{equation}
已经证明
\begin{equation}\tag{5.6.8.13}\dim E_2^{p,q}=\dim E_3^{p,q},\qquad(p,q)\neq(0,n),(2,n-1),\end{equation}
而由\Term{degree}{次数}的原因，
\begin{equation}\tag{5.6.8.14}E_3^{p,q}=E_\infty^{p,q}\qquad\text{对所有 }(p,q).\end{equation}
\SourcePage{307}
由 (5.6.8.14)，有 $b_i(X)=\dim H^i(X,\mathbf Q_\ell)=\sum_{p+q=i}\dim E_3^{p,q}$，再由 (5.6.8.13)，有
\begin{align}
b_{n-1}(X)&=\dim E_2^{0,n-1}+\dim E_2^{1,n-2}+\dim E_2^{2,n-3},\tag{5.6.8.15}\\
b_{n+1}(X)&=\dim E_2^{0,n+1}+\dim E_2^{1,n}+\dim E_3^{2,n-1}.\tag{5.6.8.16}
\end{align}
由\Term{poincare-duality}{Poincaré 对偶性}，有
\begin{equation}\tag{5.6.8.17}b_{n-1}(X)=b_{n+1}(X).\end{equation}
先注意，由 (5.6.8.8)，有 $\dim E_2^{1,n}=\dim E_2^{1,n-2}$。因此，由 (5.6.8.17)，一旦知道
\begin{align}
\dim E_2^{2,n-3}&=\dim E_2^{0,n+1},\tag{5.6.8.18}\\
\dim E_2^{0,n-1}&=\dim E_2^{2,n-1},\tag{5.6.8.19}
\end{align}
(5.6.8.12) 就得证。而这两个公式来自 (5.6.8.2)、$p$ 的\Term{smooth-fibre}{光滑纤维}的\Term{poincare-duality}{Poincaré 对偶性}，以及下面的\Term{lemma}{引理}；将该\Term{lemma}{引理}应用于 $\mathbf P^1$ 和\Term{sheaf}{层} $R^{n-3}p_*\mathbf Q_\ell$、$R^{n-1}p_*\mathbf Q_\ell$ 即可。
\end{proof}
\SourcePage{308}
\begin{SGAresult}{\Term{lemma}{引理}}{5.6.9}
（参见 SGA 5 I \S5 及 SGA 4 XVIII\SGAsourcegap{原 PDF 第308页\Term{lemma}{引理}5.6.9的“SGA 4 XVIII”之后留白，缺少具体定位。}。）设 $D$ 是\Term{algebraically-closed-field}{代数闭域} $k$ 上的一条\Term{proper}{固有}、\Term{smooth}{光滑}、\Term{connected}{连通}\Term{curve}{曲线}。记 $i:\eta\hookrightarrow D$ 为\Term{generic-point}{泛点}的\Term{inclusion}{包含映射}，$F$ 为 $D$ 上的\Term{constructible-sheaf}{可构造} $\mathbf Q_\ell$-\Term{sheaf}{层}，且 $F\xrightarrow{\sim}i_*i^*F$。则
\[H^2(D,F)\simeq H^0(D,\check F(-1))^\vee,\]
其中 $\check F\simeq i_*((i^*F)^\vee)$。
\end{SGAresult}
\begin{SGAresult}{推论}{5.6.10}
（使用 Picard--Lefschetz。）设 $p:\widetilde X\to\mathbf P^1$ 是来自一个\Term{lefschetz-pencil}{Lefschetz 铅笔}的“\Term{fibration}{纤维化}”，并假设\Term{smooth-fibre}{光滑纤维}满足条件 (LV)。则
\begin{equation}\tag{5.6.10.1}H^i(X_{\bar\eta},\mathbf Q_\ell)^\pi=\operatorname{Im}\bigl(H^i(X,\mathbf Q_\ell)\to H^i(X_{\bar\eta},\mathbf Q_\ell)\bigr).\end{equation}
\end{SGAresult}
\begin{proof}
由 5.3.4（其中使用了\Term{picard-lefschetz-formula}{Picard--Lefschetz 公式}），5.6.8 可以应用，得到
\[H^i(X_{\bar\eta},\mathbf Q_\ell)^\pi=\operatorname{Image}H^i(\widetilde X,\mathbf Q_\ell).\]
还须说明 $H^i(X,\mathbf Q_\ell)$ 与 $H^i(\widetilde X,\mathbf Q_\ell)$ 在 $H^i(X_{\bar\eta},\mathbf Q_\ell)$ 中的像相同。当 $i\leq n-1$ 时由 5.1.6，当 $i\geq n$ 时由 5.2.5，这都成立。
\end{proof}
现在进行最重要的计算：
\begin{SGAresult}{\Term{theorem}{定理}}{5.7}
在 5.6 的假设下，置
\[E^{n-2}(\Delta,\mathbf Q_\ell(j))=\text{$H^{n-2}(\Delta,\mathbf Q_\ell(j))$ 中 $H^{n-2}(X,\mathbf Q_\ell(j))$ 的像的正交补}\]
（$\Delta$ 的\Term{cohomology}{上同调}的“\Term{vanishing-cohomology}{消失部分}”，见 5.4.4），这里所取的是\Term{orthogonal-complement}{正交补}。则有\Term{direct-sum-decomposition}{直和分解}（见 4.2）
\begin{equation}\tag{5.7.1}E_2^{1,n-1}=H^1(\mathbf P^1,R^{n-1}p_*\mathbf Q_\ell)\xleftarrow[\sim]{f^*+i_*g^*}\operatorname{Prim}^n(X,\mathbf Q_\ell)\oplus E^{n-2}(\Delta,\mathbf Q_\ell(-1)).\end{equation}
\end{SGAresult}
\SourcePage{309}
\begin{proof}
由\Term{degenerate-spectral-sequence}{退化}，有
\[E_2^{1,n-1}\simeq F^1H^n(\widetilde X,\mathbf Q_\ell)/F^2H^n(\widetilde X,\mathbf Q_\ell).\]
由 (A)，有
\begin{equation}\tag{5.7.2}\begin{aligned}
F^1H^n(\widetilde X,\mathbf Q_\ell)&=\operatorname{Ker}\bigl(k^*:H^n(\widetilde X,\mathbf Q_\ell)\to H^n(X_{\bar\eta},\mathbf Q_\ell)\bigr)\\
&\overset{(5.1.2)}\simeq\operatorname{Ker}\bigl(H^n(X,\mathbf Q_\ell)\oplus H^{n-2}(\Delta,\mathbf Q_\ell(-1))\xrightarrow{m^*+h_*}H^n(X_{\bar\eta},\mathbf Q_\ell)\bigr).
\end{aligned}\end{equation}
由 (LV)，$k^*$ 又可改写为
\begin{equation}\tag{5.7.3}\begin{gathered}
\operatorname{Prim}^n(X,\mathbf Q_\ell)\oplus H^{n-2}(X,\mathbf Q_\ell(-1))\oplus H^{n-2}(\Delta,\mathbf Q_\ell(-1))\\
\xrightarrow{m^*+m^*L+h_*}H^n(X_{\bar\eta},\mathbf Q_\ell),
\end{gathered}\end{equation}
而中间的映射
\begin{equation}\tag{5.7.4}H^{n-2}(X,\mathbf Q_\ell(-1))\xrightarrow{m^*L}H^n(X_{\bar\eta},\mathbf Q_\ell)\end{equation}
是\Term{isomorphism}{同构}（见 5.2.3 及 (5.2.2.3)）。因而，记到\Term{primitive-part}{本原部分}的\Term{projection}{投影}为
\begin{equation}\tag{5.7.5}\mathfrak P:H^n(X,\mathbf Q_\ell)\longrightarrow\operatorname{Prim}^n(X,\mathbf Q_\ell),\end{equation}
就有
\begin{equation}\tag{5.7.6}F^1H^n(\widetilde X,\mathbf Q_\ell)\xrightarrow[\sim]{\mathfrak P f_*\oplus(-g_*i^*)}\operatorname{Prim}^n(X,\mathbf Q_\ell)\oplus H^{n-2}(\Delta,\mathbf Q_\ell(-1)).\end{equation}
由 (5.6.2.2) 及\Term{degenerate-spectral-sequence}{退化}，有
\begin{equation}\tag{5.7.7}F^2H^n(\widetilde X,\mathbf Q_\ell)\simeq E_2^{2,n-2}\simeq H^{n-2}(X_{\bar\eta},\mathbf Q_\ell(-1)).\end{equation}
\SourcePage{310}
另一方面，$F^2H^n(\widetilde X,\mathbf Q_\ell)=H^{n-2}(X_{\bar\eta},\mathbf Q_\ell(-1))$ 到 $H^n(\widetilde X,\mathbf Q_\ell)$ 的\Term{inclusion}{包含映射}由\Term{gysin-homomorphism}{Gysin 映射} $k_*$ 给出。（严格说来，\Term{gysin-homomorphism}{Gysin 映射}只定义为映射
\[H^{i-2}(X_{\bar\eta},\mathbf Q_\ell)\longrightarrow H^i(X\times_k k(\bar\eta),\mathbf Q_\ell),\]
但对每个\Term{algebraically-closed-extension}{代数闭扩域} $K/k$，都有\Term{canonical}{典范}\Term{isomorphism}{同构}（SGA 4 XII 5.4）
\[H^i(X,\mathbf Q_\ell)\xrightarrow{\sim}H^i(X\otimes_k K,\mathbf Q_\ell),\]
因此希望读者允许这种记号上的滥用。）根据 (5.1.3)，\Term{gysin-homomorphism}{Gysin 映射}表示为
\begin{equation}\tag{5.7.8}H^{n-2}(X_{\bar\eta},\mathbf Q_\ell(-1))\xrightarrow{m_*\oplus(-h^*)}H^n(X,\mathbf Q_\ell)\oplus H^{n-2}(\Delta,\mathbf Q_\ell(-1)).\end{equation}
由\Term{weak-lefschetz-theorem}{弱 Lefschetz 定理}，
\[m_*H^{n-2}(X_{\bar\eta},\mathbf Q_\ell(-1))=m_*m^*H^{n-2}(X,\mathbf Q_\ell(-1))=L_XH^{n-2}(X,\mathbf Q_\ell(-1))\]
是\Term{primitive-part}{本原部分}的\Term{orthogonal-complement}{正交补}，因此有\Term{commutative-diagram}{交换图}
\begin{equation}\tag{5.7.9}
\begin{tikzcd}[column sep=large]
F^2H^n(\widetilde X,\mathbf Q_\ell)\arrow[r,hook]&F^1H^n(\widetilde X,\mathbf Q_\ell)\arrow[d,"{\mathfrak P f_*\oplus(-g_*i^*)}","\sim"']\\
H^{n-2}(X_{\bar\eta},\mathbf Q_\ell(-1))\arrow[u,"k_*","\sim"']\arrow[r,"{0\oplus h^*}"']&\operatorname{Prim}^n(X,\mathbf Q_\ell)\oplus H^{n-2}(\Delta,\mathbf Q_\ell(-1)),
\end{tikzcd}
\end{equation}
这就给出 (5.7.1)，证毕。
\end{proof}
\SourcePage{311}
\textbf{注记 5.7.10。} 正如证明所示，在 5.7 中，假设 (A) 可以换成下述较弱的假设：当 $q=n$ 时，映射
\begin{equation}\tag{5.7.11}H^q(\widetilde X,\mathbf Q_\ell)\longrightarrow H^0(\mathbf P^1,R^qp_*\mathbf Q_\ell)\end{equation}
及
\begin{equation}\tag{5.7.12}H^q(\widetilde X,\mathbf Q_\ell)\longrightarrow H^q(X_{\bar\eta},\mathbf Q_\ell)\end{equation}
具有相同的核。如果不加这个条件，结论可能不成立，见 5.9。
\SGAsection{5.8}{Griffiths 同态}
\textbf{5.8.1。} 设 $f:T\to S$ 是\Term{scheme}{概形}的\Term{morphism}{态射}。\Term{leray-spectral-sequence}{Leray 谱序列}在 $H^*(T,\mathbf Q_\ell(j))$ 上定义了\Term{filtration}{滤过} $(F^i)$。定义
\begin{equation}\tag{5.8.1.1}\operatorname{Prim}^*(T/S,\mathbf Q_\ell(j))=F^1H^*(T,\mathbf Q_\ell(j)).\end{equation}
由于\Term{leray-spectral-sequence}{Leray 谱序列}位于“第一象限”，有
\begin{equation}\tag{5.8.1.2}E_\infty^{0,*}\hookrightarrow E_2^{0,*},\end{equation}
因而
\begin{equation}\tag{5.8.1.3}\operatorname{Prim}^*(T/S,\mathbf Q_\ell(j))=\operatorname{Ker}\bigl(H^*(T,\mathbf Q_\ell(j))\to H^0(S,R^*f_*\mathbf Q_\ell(j))\bigr).\end{equation}
同理，有
\begin{equation}\tag{5.8.1.4}E_\infty^{1,*}\hookrightarrow E_2^{1,*}.\end{equation}
\SourcePage{312}
复合
\begin{equation}\tag{5.8.1.5}\begin{aligned}
\operatorname{Prim}^i(T/S,\mathbf Q_\ell(j))=F^1H^i(T,\mathbf Q_\ell(j))&\longrightarrow\operatorname{gr}^1H^i(T,\mathbf Q_\ell(j))=E_\infty^{1,i-1}\\
&\lhook\joinrel\longrightarrow H^1(S,R^{i-1}f_*\mathbf Q_\ell(j))=E_2^{1,i-1}
\end{aligned}\end{equation}
记作
\begin{equation}\tag{5.8.1.6}u^i:\operatorname{Prim}(T/S,\mathbf Q_\ell(j))\longrightarrow H^1(S,R^{i-1}f_*\mathbf Q_\ell(j)).\end{equation}
\textbf{5.8.2。} 现在设 $X$ \Term{projective}{射影}、\Term{smooth}{光滑}且\Term{irreducible}{不可约}，$D=\{H_\tau\}$ 是一个\Term{hyperplane-pencil}{超平面铅笔}，其\Term{generic-member}{一般成员}及其\Term{pencil-axis}{轴}与 $X$ \Term{transverse-intersection}{横截相交}。记 $U\xhookrightarrow{\nu}\mathbf P^1$ 为\Term{projection}{投影} $p:\widetilde X\to\mathbf P^1$ 的一个\Term{nonempty}{非空}\Term{smooth-locus}{光滑开集}，令 $\widetilde X|U=\widetilde X\times_{\mathbf P^1}U$。有\Term{commutative-diagram}{交换图}
\begin{equation}\tag{5.8.2.1}
\begin{tikzcd}[column sep=large]
&\widetilde X\arrow[dl,"f"']\arrow[d,"p"]&\widetilde X|U\arrow[l,"\alpha"']\arrow[d,"p|U"]\\
X&\mathbf P^1&U\arrow[l,"\nu"]
\end{tikzcd}
\end{equation}
\SourcePage{313}
因为\Term{morphism}{态射} $p|U:\widetilde X|U\to U$ \Term{proper-morphism}{固有}且\Term{smooth-morphism}{光滑}，由\Term{cohomological-specialization}{上同调专化定理}（SGA 4 XVI 2.2），有
\begin{equation}\tag{5.8.2.2}\operatorname{Prim}^*(\widetilde X|U/U,\mathbf Q_\ell(j))=\operatorname{Ker}\bigl(H^*(\widetilde X|U,\mathbf Q_\ell(j))\to H^*(X_{\bar\eta},\mathbf Q_\ell(j))\bigr);\end{equation}
因此，根据 (5.2.3.1)，有
\begin{SGAresult}{命题}{5.8.2.3}
在 (5.8.2) 的假设下，对 $x\in H^n(X,\mathbf Q_\ell(j))$，有\Term{equivalence-relation}{等价关系}
\begin{equation}\tag{5.8.2.4}x\in\operatorname{Prim}^n(X,\mathbf Q_\ell(j))\quad\Longleftrightarrow\quad\alpha^*f^*(x)\in\operatorname{Prim}^n(\widetilde X|U/U,\mathbf Q_\ell(j)).\end{equation}
如果 $p$ 还满足 5.3 的假设 (A)，则有
\begin{equation}\tag{5.8.2.5}x\in\operatorname{Prim}^n(X,\mathbf Q_\ell(j))\quad\Longleftrightarrow\quad f^*(x)\in\operatorname{Prim}^n(\widetilde X/\mathbf P^1,\mathbf Q_\ell(j)).\end{equation}
\end{SGAresult}
\begin{proof}
(5.8.2.4) 已经证明。至于 (5.8.2.5)，假设 (A) 给出
\[H^0(\mathbf P^1,R^np_*\mathbf Q_\ell(j))\xrightarrow[\sim]{\nu^*}H^0(U,\nu^*R^np_*\mathbf Q_\ell(j)),\]
所以 [$f^*(x)$ 相对于 $\mathbf P^1$ 是\Term{primitive-part}{本原的}] 当且仅当 [$\alpha^*f^*(x)$ 相对于 $U$ 是\Term{primitive-part}{本原的}]。
\end{proof}
\begin{SGAresult}{定义}{5.8.3}
在 5.8.2 的情形下，\Term{composite-homomorphism}{复合同态}
\begin{equation}\tag{5.8.3.1}
\begin{tikzcd}[column sep=large]
\operatorname{Prim}^n(X,\mathbf Q_\ell(j))\arrow[r,"\alpha^*f^*"]\arrow[dr,"\operatorname{Grif}"']&\operatorname{Prim}^n(\widetilde X|U/U,\mathbf Q_\ell(j))\arrow[d,"u"]\\
&H^1(U,\nu^*R^{n-1}p_*\mathbf Q_\ell(j))
\end{tikzcd}
\end{equation}
称为（$X$ 相对于给定\Term{pencil}{铅笔} $D$ 的）\Term{griffiths-homomorphism}{Griffiths 同态}。
\end{SGAresult}
\SourcePage{314}
在图 (5.8.2.1) 中加入 $U$ 上\Term{fibre}{纤维}为 $X$ 的\Term{constant-family}{常值族}：
\begin{equation}\tag{5.8.3.2}
\begin{tikzcd}[column sep=large]
&\widetilde X\arrow[dl,"f"']\arrow[d,"p"]&\widetilde X|U\arrow[l,"\alpha"']\arrow[d,"p|U"]\arrow[r,"f\times p|U"]&X\times U\arrow[dl,"\operatorname{pr}_2"]\\
X&\mathbf P^1&U\arrow[l,"\nu"]&
\end{tikzcd}
\end{equation}
\begin{SGAresult}{\Term{theorem}{定理}}{5.8.4}
假设 $X$ 满足 (LV) (5.2.2)，且 $p$ 满足假设 (A) (5.3)。则对 $p$ 的每个\Term{nonempty}{非空}\Term{smooth-locus}{光滑开集} $U\xhookrightarrow{\nu}\mathbf P^1$，“模去\Term{fixed-part}{固定部分}”的\Term{griffiths-homomorphism}{Griffiths 同态}，按定义即复合
\begin{equation}\tag{5.8.4.1}\begin{gathered}
\operatorname{Prim}^n(X,\mathbf Q_\ell(j))\xrightarrow{\operatorname{Grif}}H^1(U,\nu^*R^{n-1}p_*\mathbf Q_\ell(j))\\
\longrightarrow H^1(U,E^{n-1}p_*\mathbf Q_\ell(j))\simeq\frac{H^1(U,\nu^*R^{n-1}p_*\mathbf Q_\ell(j))}{(f\times p|U)^*H^1(U,H^{n-1}(X,\mathbf Q_\ell(j))_U)},
\end{gathered}\end{equation}
是\Term{injectivity}{单射}。
\end{SGAresult}
\SourcePage{315}
\begin{proof}
由 (5.5.1)，在 $\mathbf P^1$ 上有
\begin{equation}\tag{5.8.4.2}R^{n-1}p_*\mathbf Q_\ell(j)\simeq H^{n-1}(X,\mathbf Q_\ell(j))_{\mathbf P^1}\oplus E^{n-1}p_*\mathbf Q_\ell(j),\end{equation}
在 $U$ 上有
\begin{equation}\tag{5.8.4.3}\nu^*R^{n-1}p_*\mathbf Q_\ell(j)\simeq H^{n-1}(X,\mathbf Q_\ell(j))_U\oplus\nu^*E^{n-1}p_*\mathbf Q_\ell(j),\end{equation}
所以能把 (5.8.4.1) 按所述方式改写为
\begin{equation}\tag{5.8.4.4}\begin{gathered}
\operatorname{Prim}^n(X,\mathbf Q_\ell(j))\xrightarrow{\operatorname{Grif}}H^1(U,\nu^*R^{n-1}p_*\mathbf Q_\ell(j))\\
\longrightarrow H^1(U,\nu^*E^{n-1}p_*\mathbf Q_\ell(j)).
\end{gathered}\end{equation}
由 (5.7.1)，映射（它只因假设 (A) 才有定义）
\begin{equation}\tag{5.8.4.5}u:\operatorname{Prim}^n(X,\mathbf Q_\ell(j))\longrightarrow H^1(\mathbf P^1,R^{n-1}p_*\mathbf Q_\ell)\end{equation}
是\Term{injectivity}{单射}。由于 $\mathbf P^1$ \Term{simply-connected}{单连通}，\Term{decomposition}{分解} (5.8.4.2) 给出
\begin{equation}\tag{5.8.4.6}H^1(\mathbf P^1,R^{n-1}p_*\mathbf Q_\ell)\simeq H^1(\mathbf P^1,E^{n-1}p_*\mathbf Q_\ell(j)),\end{equation}
从而可以把 (5.8.4.4) 补成\Term{commutative-diagram}{交换图}
\SourcePage{316}
\begin{equation}\tag{5.8.4.7}
\begin{tikzcd}[column sep=large]
\operatorname{Prim}^n(X,\mathbf Q_\ell(j))\arrow[r,"\operatorname{Grif}"]\arrow[d,"u"',hook]&H^1(U,\nu^*R^{n-1}p_*\mathbf Q_\ell(j))\arrow[d]\\
H^1(\mathbf P^1,E^{n-1}p_*\mathbf Q_\ell(j))\arrow[r,"\nu^*"']&H^1(U,\nu^*E^{n-1}p_*\mathbf Q_\ell(j)).
\end{tikzcd}
\end{equation}
现在只需说明下面的映射为\Term{injectivity}{单射}。这由
\begin{equation}\tag{5.8.4.8}E^{n-1}p_*\mathbf Q_\ell(j)\simeq\nu_*\nu^*E^{n-1}p_*\mathbf Q_\ell(j)\end{equation}
及以下（众所周知的！）\Term{lemma}{引理}\Term{pushout}{推出}。
\end{proof}
\begin{SGAresult}{\Term{lemma}{引理}}{5.8.5}
设 $\nu:A\to B$ 是\Term{topos}{拓扑斯}的\Term{morphism}{态射}，$F$ 是 $A$ 上的一个\Term{sheaf}{层}。则\Term{canonical-map}{典范映射}
\begin{equation}\tag{5.8.5.1}H^1(B,\nu_*F)\longrightarrow H^1(A,F)\end{equation}
为\Term{injectivity}{单射}。
\end{SGAresult}
\begin{proof}
在第一象限的\Term{spectral-sequence}{谱序列}中，有
\begin{equation}\tag{5.8.5.2}E_2^{1,0}\simeq E_\infty^{1,0}.\end{equation}
因此对\Term{leray-spectral-sequence}{Leray 谱序列}
\[E_2^{p,q}=H^p(B,R^q\nu_*F)\Longrightarrow H^{p+q}(A,F),\]
\SourcePage{317}
有
\begin{equation}\tag{5.8.5.3}E_2^{1,0}=H^1(B,\nu_*F)\simeq E_\infty^{1,0}=F^1H^1(A,F).\end{equation}
取 $F=\nu^*E$，其中 $E$ 是 $B$ 上的一个\Term{sheaf}{层}，得到
\begin{equation}\tag{5.8.5.4}H^1(B,\nu_*\nu^*E)\hookrightarrow H^1(A,\nu^*E).\end{equation}
\end{proof}
\textbf{5.8.6。} 设 $L/k(\eta)$ 为 $\mathbf P^1$ 的\Term{rational-function-field}{有理函数域}的一个\Term{finite-extension}{有限扩张}，置
\begin{equation}\tag{5.8.6.1}S=\text{$\mathbf P^1$ 在 $L$ 中的正规化}.\end{equation}
在 (5.8.3.2) 的情形中作\Term{base-change}{基变换} $S\to\mathbf P^1$：
\begin{equation}\tag{5.8.6.1 bis}
\begin{tikzcd}[column sep=large]
&\widetilde X\arrow[dl,"f"']\arrow[d,"p"]&\widetilde X_S\arrow[l,"\alpha_S"']\arrow[d,"p_S"]\\
X&\mathbf P^1&S\arrow[l,"\pi"]
\end{tikzcd}
\end{equation}
对 $S$ 的每个\Term{open-subset}{开集} $V\xhookrightarrow{\mu}S$，使 $p_S$ 在其上\Term{smooth-morphism}{光滑}，在 (5.8.3.2) 中作\Term{base-change}{基变换} $V\xrightarrow{\pi\mu}\mathbf P^1$，得到
\begin{equation}\tag{5.8.6.2}
\begin{tikzcd}[column sep=large]
&\widetilde X\arrow[dl,"f"']\arrow[d,"p"]&\widetilde X_V\arrow[l,"\alpha_V"']\arrow[d,"p_V"]\arrow[r,"f\times p_V"]&X\times V\arrow[dl,"\operatorname{pr}_2"]\\
X&\mathbf P^1&V\arrow[l,"\pi\mu"]&
\end{tikzcd}
\end{equation}
\SourcePage{318}
且仍有（见 (5.8.2.3)）
\begin{equation}\tag{5.8.6.3}(\alpha_Vf)^*\operatorname{Prim}^n(X,\mathbf Q_\ell(j))\subset\operatorname{Prim}^n(\widetilde X_V/V,\mathbf Q_\ell(j)),\end{equation}
因此仍可谈论\Term{griffiths-homomorphism}{Griffiths 同态}
\begin{equation}\tag{5.8.6.4}\operatorname{Prim}^n(X,\mathbf Q_\ell(j))\xrightarrow{\operatorname{Grif}}H^1(V,R^{n-1}p_{V*}\mathbf Q_\ell(j)).\end{equation}
\begin{SGAresult}{推论}{5.8.7}
沿用上述记号，并在 5.8.4 的假设下，“模去\Term{fixed-part}{固定部分}”的\Term{griffiths-homomorphism}{Griffiths 同态}，即复合
\begin{equation}\tag{5.8.7.1}\begin{gathered}
\operatorname{Prim}^n(X,\mathbf Q_\ell(j))\xrightarrow{\operatorname{Grif}}H^1(V,R^{n-1}p_{V*}\mathbf Q_\ell(j))\\
\longrightarrow H^1(V,(\pi\mu)^*E^{n-1}p_*\mathbf Q_\ell(j))\simeq\frac{H^1(V,R^{n-1}p_{V*}\mathbf Q_\ell(j))}{(f\times p_V)^*H^1(V,H^{n-1}(X,\mathbf Q_\ell(j))_V)},
\end{gathered}\end{equation}
是\Term{injectivity}{单射}。
\end{SGAresult}
\begin{proof}
由\Term{cohomological-specialization}{上同调专化定理}，有
\[R^{n-1}p_{V*}\mathbf Q_\ell(j)=(\pi\mu)^*R^{n-1}p_*\mathbf Q_\ell(j),\]
\SourcePage{319}
所以由\Term{decomposition}{分解} (5.8.4.2)，(5.8.7.1) 确可如上改写为
\[\begin{gathered}\operatorname{Prim}^n(X,\mathbf Q_\ell(j))\xrightarrow{\operatorname{Grif}}H^1(V,(\pi\mu)^*R^{n-1}p_*\mathbf Q_\ell(j))\\
\longrightarrow H^1(V,(\pi\mu)^*E^{n-1}p_*\mathbf Q_\ell(j)).\end{gathered}\]
因此只须对任意小的 $V\xhookrightarrow{\mu}S$，证明
\[\operatorname{Prim}^n(X,\mathbf Q_\ell(j))\longrightarrow H^1(V,\mu^*\pi^*E^{n-1}p_*\mathbf Q_\ell(j))\]
的\Term{injectivity}{单射}性。可归结到下述情形：$U\xhookrightarrow{\nu}\mathbf P^1$ 是 $p$ 的一个\Term{smooth-locus}{光滑开集}，$V\xrightarrow{\pi_U}U$ 是 $U$ 在 $L$ 中的\Term{normalization}{正规化}。由\Term{cartesian-diagram}{笛卡尔}\Term{commutative-diagram}{交换图}
\[
\begin{tikzcd}
V\arrow[r,hook,"\mu"]\arrow[d,"\pi_U"']&S\arrow[d,"\pi"]\\
U\arrow[r,hook,"\nu"']&\mathbf P^1
\end{tikzcd}
\]
\SourcePage{320}
得出
\[\mu^*\pi^*E^{n-1}p_*\mathbf Q_\ell(j)=\pi_U^*\nu^*E^{n-1}p_*\mathbf Q_\ell,\]
所以 (5.8.7.3) \Term{decomposition}{分解}为
\[\begin{gathered}\operatorname{Prim}^n(X,\mathbf Q_\ell(j))\longrightarrow H^1(U,\nu^*E^{n-1}p_*\mathbf Q_\ell(j))\\
\xrightarrow{\pi_U^*}H^1(V,\pi_U^*\nu^*E^{n-1}p_*\mathbf Q_\ell(j)).\end{gathered}\]
由 (5.8.4)，只须验证
\[\pi_U^*:H^1(U,\nu^*E^{n-1}p_*\mathbf Q_\ell(j))\longrightarrow H^1(V,\pi_U^*\nu^*E^{n-1}p_*\mathbf Q_\ell(j))\]
为\Term{injectivity}{单射}。因为 $\pi_U:V\to U$ \Term{finite-morphism}{有限}且\Term{flat}{平坦}，\Term{trace-theory}{迹理论} [SGA 4 XVII 6.2] 对 $U$ 上的每个 $\mathbf Q_\ell$-\Term{sheaf}{层} $F$（例如 $F=\nu^*E^{n-1}p_*\mathbf Q_\ell(j)$）给出一个映射
\[\operatorname{tr}_{\pi_U}:H^1(V,\pi_U^*F)\longrightarrow H^1(U,F(j)),\]
满足
\[\operatorname{tr}_{\pi_U}\circ\pi_U^*=\deg(\pi_U)\cdot\operatorname{id},\]
这就表明 $H^1(U,F)\to H^1(V,\pi_U^*F)$ 是\Term{injectivity}{单射}，证明完成。
\end{proof}
\SourcePage{321}
\SGAsection{5.9}{假设 (A) 不成立的一个例子}
设 $X$ 是 $\mathbf P^3$ 中的一个\Term{smooth-quadric-surface}{光滑二次曲面}。则 $X$ 满足 (LV)（5.2.4）。取一个\Term{hyperplane-pencil}{超平面铅笔}（愿意的话，可以取为\Term{lefschetz-pencil}{Lefschetz 铅笔}），其\Term{generic-member}{一般成员}及其\Term{pencil-axis}{轴}与 $X$ \Term{transverse-intersection}{横截相交}。那么（如 Deligne 向我指出的），\Term{projection}{投影} $p:\widetilde X\to\mathbf P^1$ 不满足 (A)。事实上，$R^2p_*\mathbf Q_\ell$ 并不\Term{constant-sheaf}{常值}，因为\Term{singular-fibre}{奇异纤维}都是\Term{singular-conic}{奇异二次曲线}，所以\Term{reducible}{可约}；对于\Term{lefschetz-pencil}{Lefschetz 铅笔}，每条这样的\Term{fibre}{纤维}都是两条直线的并，它们在一点相交（在这个简单情形中，这就是\Term{lefschetz-pencil}{Lefschetz 铅笔}的定义）。因而
\[\dim_{\mathbf Q_\ell}H^2(\text{奇异纤维},\mathbf Q_\ell)=2,\]
尽管
\[\dim_{\mathbf Q_\ell}H^2(\text{光滑纤维},\mathbf Q_\ell)=1.\]
此外，\Term{theorem}{定理} 5.8.4 的结论在这里也不成立。记 $U\xhookrightarrow{\nu}\mathbf P^1$ 为 $p$ 的一个\Term{nonempty}{非空}\Term{smooth-locus}{光滑开集}。因为 $p$ 在 $U$ 上的\Term{fibre}{纤维}是\Term{genus-zero}{亏格零}的\Term{smooth-curve}{光滑曲线}，有
\[\nu^*R^1p_*\mathbf Q_\ell=0,\]
从而
\[H^1(U,\nu^*R^1p_*\mathbf Q_\ell)=0;\]
\SourcePage{322}
因此\Term{griffiths-homomorphism}{Griffiths 同态}
\[\operatorname{Prim}^2(X,\mathbf Q_\ell)\xrightarrow{\operatorname{Grif}}H^1(U,\nu^*R^1p_*\mathbf Q_\ell)\]
为零。但在\Term{quadric}{二次曲面}上，有
\[\dim_{\mathbf Q_\ell}\operatorname{Prim}^2(X,\mathbf Q_\ell)=1.\]
同样可见，5.7 的结论在这里也不成立。
